\section{РАЗРАБОТКА СИСТЕМЫ ИНДЕКСАЦИИ БЛОКЧЕЙН ДАННЫХ}

\subsection{Реализация серверной части системы}

Практическая реализация проектируемой системы выполнена в виде серверного
приложения на языке Rust. Серверная часть объединяет загрузку конфигурации,
инициализацию подключений к базе данных и RPC-источнику, запуск HTTP API,
фоновых процессов обработки и средств наблюдаемости. Такой способ
организации соответствует выбранной архитектуре модульного монолита со
stateless-характером исполнения и обеспечивает единый жизненный цикл всех
компонентов приложения.

Исходный код программной реализации размещен в открытом репозитории
Bitcoin Blockchain Indexer \cite{project_github}. Репозиторий используется как
основной артефакт фиксации реализации backend-приложения, миграций,
конфигурационных файлов, вспомогательных средств запуска и материалов,
связанных с проектированием системы.

Ключевая роль серверной части заключается в координации нескольких
взаимосвязанных процессов:
\begin{enumerate}
    \item приема и валидации конфигурации запуска;
    \item синхронизации заданий индексирования из YAML-конфигурации и управляющего API;
    \item выполнения индексации подтвержденных блоков и транзакций;
    \item отдельной синхронизации состояния mempool;
    \item публикации API, метрик и диагностической информации.
\end{enumerate}

С инженерной точки зрения такое построение позволяет рассматривать систему как
единый прикладной контур обработки данных, в котором критическое состояние
сохраняется во внешней базе данных, а экземпляр приложения может быть
перезапущен без потери консистентности. Тем самым реализация серверной части
непосредственно поддерживает сформулированные ранее требования к
восстанавливаемости и управляемости системы.

Ключевые инженерные решения, использованные при реализации серверной части,
приведены в таблице~\ref{tab:implementation-decisions}. В ней архитектурные
требования соотносятся с конкретными механизмами кода и базы данных.

\begin{footnotesize}
\setlength{\tabcolsep}{3pt}
\begin{longtable}{|p{0.23\textwidth}|p{0.27\textwidth}|p{0.23\textwidth}|p{0.15\textwidth}|}
    \caption{Ключевые инженерные решения реализации}\label{tab:implementation-decisions}\\ \hline
    Решение & Назначение & Реализация & Проверка \\ \hline
    \endfirsthead
    \caption*{Продолжение таблицы \ref{tab:implementation-decisions}}\\ \hline
    Решение & Назначение & Реализация & Проверка \\ \hline
    \endhead
    \hline
    \endfoot
    \hline
    \endlastfoot
    Advisory lock на состояние цепи &
    исключить параллельную запись конфликтующих высот &
    advisory lock PostgreSQL в транзакции записи блока &
    pipeline-тесты \\ \hline
    Идемпотентная запись &
    безопасно повторять обработку блока или job &
    \texttt{ON CONFLICT}, первичные и уникальные ключи &
    storage-тесты \\ \hline
    Разделение статусов транзакций &
    не смешивать confirmed, mempool, dropped и orphaned данные &
    поле \texttt{status} и отдельные API-фильтры &
    data API тесты \\ \hline
    Пересборка агрегатов после reorg &
    восстановить UTXO и балансы по канонической ветке &
    пометка orphaned и replay canonical blocks &
    reorg-тесты \\ \hline
    Прогресс jobs в БД &
    сохранить управляемость после рестарта приложения &
    поле \texttt{progress\_height} и переходы job &
    jobs API тесты \\ \hline
    Метрики Prometheus &
    наблюдать лаг, ошибки и объем обработки &
    \texttt{/metrics}, счетчики и gauge-показатели &
    smoke-проверка \\ \hline
\end{longtable}
\end{footnotesize}

\subsection{Реализация механизма индексации блоков и транзакций}

Основной процесс обработки подтвержденных данных реализован в модуле
индексатора.
Он получает данные о блоках от Bitcoin Core, выполняет декомпозицию транзакций,
входов и выходов, после чего сохраняет результат в PostgreSQL в рамках одной
транзакции. Такой подход обеспечивает атомарность записи и предотвращает
частичное появление несогласованных данных.

Алгоритм пакетной индексации подтвержденных данных представлен на
рисунке~\ref{alg:index-batch}. В нем отражены ключевые стадии: определение
диапазона высот, проверка расхождения цепи, последовательная запись блоков и
обновление прогресса задания индексирования.

\begin{figure}
\begin{minipage}{\linewidth}
\begin{algorithm}[H]
    \SetAlgoVlined
    \KwData{конфигурация job, состояние БД, RPC-клиент}
    \KwResult{обновленные данные подтвержденной цепи и прогресс задания}
    получить текущую высоту узла\;
    определить диапазон высот для следующего батча\;
    проверить расхождение цепи в окне \texttt{reorg\_depth}\;
    \If{обнаружен reorg}{
        пометить неканонические блоки и транзакции\;
        пересобрать агрегаты UTXO и балансов\;
        откатить \texttt{progress\_height}\;
    }
    \ForEach{высота в выбранном диапазоне}{
        получить блок по RPC\;
        декомпозировать транзакции, входы и выходы\;
        сохранить блок и связанные сущности в одной транзакции БД\;
        обновить UTXO и баланс адресов\;
        обновить \texttt{progress\_height}\;
    }
\end{algorithm}
\end{minipage}
\caption{Алгоритм пакетной индексации данных подтвержденной цепи}
\label{alg:index-batch}
\end{figure}

Особое внимание в реализации уделено идемпотентности. Повторная обработка уже
сохраненной высоты не должна приводить к появлению дубликатов или искажению
адресных агрегатов. Для этого используются ограничения целостности базы
данных, транзакционная запись и согласованное обновление производных
сущностей. Таким образом, реализация индексатора обеспечивает не только
накопление данных, но и воспроизводимость результата обработки.

\subsection{Реализация управления заданиями индексирования, mempool и reorg}

Управление заданиями индексирования реализовано как отдельный функциональный
контур, обеспечивающий хранение конфигурации заданий, контроль их состояния и периодический запуск
батчей индексации. Каждое задание имеет собственный идентификатор, режим
работы, текущий статус, информацию о прогрессе и диагностические данные о
последней ошибке. Это позволяет использовать задание как единицу управляемости и
наблюдаемости процесса индексации.

Наряду с обработкой подтвержденной цепи в системе реализован отдельный контур
синхронизации mempool. Его задача состоит в периодическом опросе узла,
сохранении неподтвержденных транзакций и пометке исчезнувших записей как
\texttt{dropped}. При этом агрегаты подтвержденной цепи не смешиваются с mempool,
поскольку их семантика принципиально различается.

Алгоритм синхронизации mempool представлен на
рисунке~\ref{alg:mempool-sync}.

\begin{figure}
\begin{minipage}{\linewidth}
\begin{algorithm}[H]
    \SetAlgoVlined
    \KwData{список известных mempool tx, ответ \texttt{getrawmempool}}
    \KwResult{актуальное состояние mempool в БД}
    получить текущий список \texttt{txid} из mempool\;
    определить новые и исчезнувшие транзакции\;
    \ForEach{новый \texttt{txid}}{
        запросить verbose-представление транзакции\;
        сохранить запись со статусом \texttt{mempool}\;
        сохранить входы и выходы для будущей фильтрации по адресу\;
    }
    \ForEach{исчезнувший \texttt{txid}}{
        пометить запись как \texttt{dropped}, если она не подтверждена\;
    }
\end{algorithm}
\end{minipage}
\caption{Алгоритм синхронизации mempool}
\label{alg:mempool-sync}
\end{figure}

Для обеспечения корректности данных канонической цепи в системе реализована обработка
reorg. При обнаружении расхождения между сохраненной в БД цепью и актуальным
состоянием узла неканонические блоки переводятся в статус \texttt{orphaned},
затронутым транзакциям присваивается соответствующая маркировка, а производные
агрегаты пересобираются по оставшейся канонической ветке. Одновременно
прогресс задания индексирования откатывается до согласованной высоты, после чего индексация
продолжается уже по новой ветви цепи.

\subsection{Реализация API, конфигурирования и мониторинга}

Внешнее взаимодействие с системой реализовано через REST API, которое
предоставляет операции управления заданиями индексирования, диагностики узлов и получения
проиндексированных данных. Для сопровождения системы предусмотрены также
административная панель и CLI, работающие поверх того же HTTP-контракта.

Конфигурирование серверного приложения выполняется через YAML-файл и переменные окружения.
В конфигурации задаются параметры API, настройки подключения к RPC-узлу,
ограничения параллелизма, параметры заданий индексирования и режимы обработки mempool и reorg.
Такой подход делает конфигурацию внешней по отношению к приложению и хорошо
согласуется с требованиями контейнерного развертывания.

Наблюдаемость системы обеспечивается за счет JSON-логирования и публикации
метрик Prometheus. В частности, экспортируются показатели текущей высоты цепи,
прогресса заданий индексирования, лага индексации, числа обработанных блоков и транзакций,
латентности RPC-вызовов и ошибок ключевых подсистем. Это позволяет использовать
разработанную систему как основу для реального эксплуатационного контура и
подтверждает соответствие реализованных средств сопровождения исходным
эксплуатационным требованиям.
